\begin{quote}
\textbf{Maturity: {[}Implemented{]} mechanism, {[}Hypothesis{]}
outcomes.} This paper gives the organizational-design rationale behind
Paper II. In February 2026 the mechanism it motivates was specified; it
is now built, the mentor-mentee loop, adapter generation, and the
on-chain registry all cite source and deployed contracts. What remains
untested is whether the \emph{organizational-learning outcomes} the
theory predicts actually emerge, and those stay explicit hypotheses with
a stated methodology.
\end{quote}

\subsection{Abstract}\label{abstract}

Federated learning treats aggregation as a stateless operation,
discarding the organizational intelligence that human institutions
encode through mentorship, hierarchical knowledge transfer, and
double-loop learning. This paper provides the organizational-design
rationale for the Paraconsistent Consensus protocol of Paper II, arguing
that three pillars of organizational learning theory, Senge's systems
thinking, Nonaka and Takeuchi's SECI knowledge-creation spiral, and
Argyris and Schön's double-loop learning, map to concrete architectural
patterns for decentralized agent swarms. The central contribution is the
Central Oracle: a knowledge aggregation hub, first proposed in the
author's 2023 internal note, that maps to BFT finality checkpoints in
the Citrate Network. The Oracle maintains performance profiles,
generates targeted LoRA-adapter mentorship signals, and improves its own
aggregation strategy over time, realizing double-loop learning at the
protocol level. We are explicit about the boundaries of the analogy: the
SECI mappings are structural analogies that motivated the architecture,
not functional equivalences that validate it. Unlike the February 2026
draft, the mechanism is now implemented, we cite the adapter and mentor
code and the deployed contracts, but the organizational-outcome
predictions remain hypotheses, and we report no simulation results in
their place.

\textbf{Keywords:} organizational learning, mentorship, agent swarms,
federated meta-learning, Central Oracle, knowledge aggregation,
double-loop learning, SECI model, Citrate Network

\subsection{1. Introduction}\label{introduction}

The aggregation server in Federated Averaging {[}2{]} computes a
weighted mean and forgets. It maintains no model of which clients excel
at which tasks, provides no targeted guidance, and cannot improve its
own aggregation strategy. This is not a limitation of the mathematics;
it is a limitation of the organizational design.

Human organizations solved this through mentorship: experienced
practitioners observe individual strengths, provide targeted guidance,
and improve as teachers through the act of teaching {[}3, 4{]}. The
recursive relationship between teaching quality and student outcomes is
the engine that drives organizations from competent to excellent
{[}5{]}. This paper asks: what organizational-design principles, drawn
from decades of management science, should govern the architecture of a
decentralized learning system?

We answer by formalizing the Mentorship Protocol, which operationalizes
three pillars of organizational learning theory as agent-swarm patterns.
From Senge {[}3{]}, organizational intelligence emerges from interaction
patterns between individuals, not from any individual's capability
alone. From Nonaka and Takeuchi {[}4{]}, the SECI spiral transforms
tacit expertise into explicit, transferable knowledge. From Argyris and
Schön {[}5{]}, the distinction between single-loop learning (correcting
errors within a framework) and double-loop learning (questioning the
framework itself).

The central contribution is the Central Oracle, a knowledge aggregation
hub that collects observations from all agents, maintains performance
profiles, and generates targeted mentorship signals. In the Citrate
Network this Oracle is realized as the BFT finality checkpoint: a
distributed, trustless commitment that periodically crystallizes the
network's collective intelligence and propagates LoRA-adapter mentorship
signals to individual node-models. The concept was first proposed in the
author's 2023 internal note and is updated here to reflect its
implementation path through the consensus architecture (Paper I) and
paraconsistent aggregation (Paper II).

\textbf{Prior work.} da Silva {[}6{]} formalized Senge's Fifth
Discipline from a multi-agent-systems perspective (JASSS), using the
SMART framework and Z specification. da Silva's findings, that
learning-organization agents must be honest, cooperative, and tenacious,
and that trust is fundamental, align with the Byzantine-fault-tolerance
requirements of our checkpoint-based system. We extend that
formalization to federated learning, where the learning organization is
a decentralized network and the knowledge-transfer mechanism is
LoRA-adapter generation at consensus checkpoints.

\textbf{Implementation status.} This paper is design rationale, the
\emph{why} behind the \emph{what} of Paper II. Sections marked
\textbf{{[}Rationale{]}} provide motivation; \textbf{{[}Hypothesis{]}}
flags testable predictions; \textbf{{[}Implemented{]}} cites code or a
deployed contract. The consensus infrastructure is implemented (Paper
I), and, unlike the February draft, the learning mechanism is now built
too (Paper II); the organizational-outcome predictions remain untested.

\subsection{2. Organizational Learning
Foundations}\label{organizational-learning-foundations}

\subsubsection{2.1 Senge's Five
Disciplines}\label{senges-five-disciplines}

\textbf{{[}Rationale{]}} Senge's \emph{The Fifth Discipline} {[}3{]}
identifies five practices that distinguish learning organizations:
personal mastery, mental models, shared vision, team learning, and
systems thinking. The fifth is integrative: organizations fail not
because individuals are incompetent but because interactions between
individuals produce emergent dysfunctions no single participant can
observe or correct.

For agent swarms this applies directly. A swarm of individually capable
models can produce collectively poor outcomes if contributions interfere
destructively, the aggregation problem that motivates paraconsistent
consensus (Paper II). Systems thinking demands that the orchestration
layer model not just individual performance but the interaction patterns
between agents: which complement each other, where redundancy helps, and
where contradiction signals genuine disagreement versus noise. da
Silva's formalization {[}6{]} shows that systems thinking requires
agents that reason about organizational structure, not just local state,
and that this depends on honest, cooperative behavior that emerges only
in trust-rich environments. In Citrate, the BFT consensus mechanism
provides the trust substrate: only nodes with blue scores above a
threshold have demonstrated honest participation, and only these can
serve as mentors (Section 4.2). This is da Silva's trust requirement
realized through consensus.

\subsubsection{2.2 The SECI Knowledge-Creation
Spiral}\label{the-seci-knowledge-creation-spiral}

\textbf{{[}Rationale{]}} Nonaka and Takeuchi {[}4{]} formalize knowledge
creation as a spiral through four modes: Socialization (tacit → tacit),
Externalization (tacit → explicit), Combination (explicit → explicit),
and Internalization (explicit → tacit).

\textbf{Table 1. SECI model → agent-swarm mapping.}

{\def\LTcaptype{none} % do not increment counter
\begin{longtable}[]{@{}
  >{\raggedright\arraybackslash}p{(\linewidth - 6\tabcolsep) * \real{0.1486}}
  >{\raggedright\arraybackslash}p{(\linewidth - 6\tabcolsep) * \real{0.2973}}
  >{\raggedright\arraybackslash}p{(\linewidth - 6\tabcolsep) * \real{0.3108}}
  >{\raggedright\arraybackslash}p{(\linewidth - 6\tabcolsep) * \real{0.2432}}@{}}
\toprule\noalign{}
\begin{minipage}[b]{\linewidth}\raggedright
SECI mode
\end{minipage} & \begin{minipage}[b]{\linewidth}\raggedright
Knowledge transition
\end{minipage} & \begin{minipage}[b]{\linewidth}\raggedright
Agent-swarm mechanism
\end{minipage} & \begin{minipage}[b]{\linewidth}\raggedright
Analogy strength
\end{minipage} \\
\midrule\noalign{}
\endhead
\bottomrule\noalign{}
\endlastfoot
Socialization & Tacit → Tacit & Agents observe peer embeddings via DAG
gossip & Weak (see note) \\
Externalization & Tacit → Explicit & Internal model state encoded as
embedding in a block & Moderate \\
Combination & Explicit → Explicit & Paraconsistent aggregation at BFT
checkpoint combines embeddings & Strong \\
Internalization & Explicit → Tacit & Node applies a LoRA adapter,
modifying internal behavior & Strong \\
\end{longtable}
}

\textbf{Note on analogy strength.} We rate each mapping because SECI
describes human knowledge processes, not computational ones.
Socialization is the weakest: in Nonaka's framework it involves shared
physical experience, apprentices learning by observing masters, whereas
gossip-protocol propagation is data broadcast, not shared experience. We
include it for completeness without claiming it captures the richness of
human socialization. Combination is the strongest: Nonaka's synthesis of
explicit knowledge from multiple sources into new explicit knowledge is
precisely what paraconsistent aggregation does when it combines node
embeddings into routing weights and Belnap state vectors at checkpoints.
Internalization is also strong: applying a LoRA adapter transforms
explicit, transferable knowledge (the adapter weights) into modified
internal behavior, Nonaka's ``learning by doing.''

\subsubsection{2.3 Single-Loop and Double-Loop
Learning}\label{single-loop-and-double-loop-learning}

\textbf{{[}Rationale{]}} Argyris and Schön {[}5{]} distinguish two
learning modes. Single-loop learning corrects errors within a framework:
the model predicted incorrectly, so we adjust weights. Double-loop
learning questions the framework: the model keeps failing on this input
class, so perhaps the routing strategy is wrong or the node should
specialize. In the Mentorship Protocol, single-loop learning corresponds
to LoRA-adapter generation, a targeted correction for a specific
weakness, and double-loop learning corresponds to the meta-model
modifying its own routing weights, changing which nodes receive which
queries. Each checkpoint is an opportunity for both. This is the
organizational argument for why the routing model must be trainable
rather than static: a fixed routing function can only perform
single-loop corrections, while a trainable one can perform the
double-loop restructuring Argyris and Schön identify as the primary
driver of organizational transformation.

\subsection{3. The Central Oracle}\label{the-central-oracle}

\subsubsection{3.1 Architecture and
History}\label{architecture-and-history}

The Central Oracle of Truth and Knowledge was first proposed in the
author's 2023 internal note as the organizational hub of a
mentorship-driven swarm, maintaining three structures: a knowledge base
aggregating observations from all agents, a performance profile tracking
each agent's accuracy across input classes, and a mentorship registry
mapping mentor-mentee pairs by complementary strengths.

\textbf{{[}Implemented{]}} In the Citrate Network the Oracle is not a
server process; it is the BFT finality checkpoint itself, and its three
structures now have concrete homes. The knowledge base is the committed
set of embeddings; the performance profile and macro-phase are fields of
the \texttt{LearningCheckpoint}
(\texttt{core/learning/src/checkpoint.rs:14}, with the per-node metrics
root at \texttt{:45}); and the mentorship registry is the on-chain
adapter registry in \textbf{\texttt{LoRAFactory.sol}} (deployed
\texttt{0x6e564d22…}), not the ``0x1003 precompile'' the February draft
named, which held \texttt{adapterModelCommitment} and
\texttt{verifyAdapterAt} for proof-backed adapter commitment. Pool
membership and cycles are managed by \texttt{LearningPool.sol} (deployed
\texttt{0xfc514b82…}), and mentor-mentee pairing is matched on-chain by
\texttt{MentorMatcher.sol}. The Oracle is not a single point of failure:
a checkpoint requires 67 of 100 validator signatures (Paper I §2.3),
making it trustless and verifiable.

\textbf{Honesty note on ``checkpoint.''} \texttt{LearningPool} exposes
\textbf{cycles} (\texttt{startCycle}/\texttt{endCycle}), not literal
per-block checkpoints; the proof-backed commitment semantics the Oracle
relies on live in \texttt{LoRAFactory}'s
\texttt{adapterModelCommitment}/\texttt{verifyAdapterAt}, and the
finality checkpoint that binds them is the BFT committee checkpoint of
Paper I, whose default interval is 50 blocks. Where this paper says ``at
each checkpoint,'' read it as ``at each learning cycle boundary,
committed against a BFT checkpoint,'' not as a per-block event.

\subsubsection{3.2 From Centralized Oracle to Distributed
Checkpoint}\label{from-centralized-oracle-to-distributed-checkpoint}

The migration from the 2023 centralized Oracle to the 2026 distributed
checkpoint follows a pattern familiar in organizational design: a
function that begins as a single role (chief knowledge officer, lead
mentor, architect) is distributed across the organization as processes
mature. In the centralized formulation the Oracle was a single server
aggregating client updates; in the distributed formulation its functions
are performed collectively by the finality committee, with BFT consensus
ensuring no single member can corrupt the aggregation. This addresses
the primary criticism of the original concept, that it reintroduces the
single point of failure decentralization was meant to eliminate, by
providing the Oracle's aggregation and mentorship-generation functions
while inheriting the consensus safety guarantees (Paper I §6.1).

The organizational reading matters because it explains a design choice
that would otherwise look arbitrary. A naive decentralization would
simply delete the Oracle and let each node train in isolation, which is
what plain FedAvg approximates when its server is a thin averaging step.
But an organization that deletes its knowledge-coordination function
does not become more capable; it becomes a collection of individuals who
cannot learn from each other. The distributed checkpoint keeps the
coordination function while removing the trusted custodian, the same
move a mature institution makes when it turns a founder's tacit judgment
into an explicit, auditable process that survives the founder. The BFT
committee is that process: it performs the Oracle's role without any
member being the Oracle, so the coordination benefit of centralization
is retained while its failure mode is not.

\subsubsection{3.3 Mentor-Mentee
Assignment}\label{mentor-mentee-assignment}

\textbf{{[}Implemented{]}} The Oracle assigns mentor-mentee
relationships by complementary performance profiles. Let Pᵢ(c) denote
node i's accuracy on input class c and P̄(c) the network median. A node
is a candidate mentor for c if Pᵢ(c) \textgreater{} P̄(c) + δ, and a
candidate mentee if Pᵢ(c) \textless{} P̄(c) − δ. The mentorship signal is
a LoRA adapter generated from the gradient of the collective loss
restricted to class c, compressed to rank r (default 16 {[}9{]}). This
is realized in \texttt{core/learning/src/mentor.rs:3}, which at each
checkpoint identifies mentor-mentee pairs by a complementarity metric,
\texttt{\textbar{}shared\_domains\textbar{}\ ·\ (mentor.accuracy\ −\ mentee.accuracy)},
and forms a delta adapter as the element-wise difference of mentor and
mentee embeddings. Assignment is dynamic: pairs are re-evaluated each
cycle, and a node that was a mentee may become a mentor after receiving
and applying adapters. This mobility is the organizational-learning
property: the system rewards improvement, not just initial capability,
mapping to Senge's personal mastery {[}3{]}.

\subsection{4. Security and Trust}\label{security-and-trust}

\subsubsection{4.1 The Mentorship Trust
Surface}\label{the-mentorship-trust-surface}

The mentor-mentee dynamic introduces a trust surface flat aggregation
avoids: a malicious mentor could generate poisoned adapters that
systematically degrade mentee performance. This is not theoretical, the
federated-learning literature documents data- and model-poisoning
attacks where Byzantine participants manipulate aggregated updates
{[}10{]}.

\subsubsection{4.2 Trust Mechanisms}\label{trust-mechanisms}

\textbf{{[}Implemented{]}} The protocol addresses this through three
mechanisms grounded in the Paper I infrastructure. \textbf{Blue-score
gating:} only nodes with blue scores above a configurable threshold can
serve as mentor candidates, a consensus-derived trust metric matching da
Silva's honesty requirement {[}6{]}. \textbf{Adapter verification:} all
adapters are committed on-chain with deterministic hashes and can be
re-derived from the committed embeddings and routing state; any node can
verify by replaying generation, and the verify-and-rollback path is
implemented (\texttt{core/learning/src/adapters.rs}:
\texttt{apply\_lora}, \texttt{remove\_lora},
\texttt{verify\_lora\_hash}), with the optimistic fraud-proof window of
Paper I §3.3 for disputes and on-chain challenge in
\texttt{AggregationChallenge.sol}. \textbf{Performance-regression
detection:} the routing model tracks per-node performance before and
after adapter application; systematic regression triggers adapter
revocation, the organizational analogue of terminating mentorship that
consistently produces worse outcomes.

\subsubsection{4.3 Byzantine Mentorship}\label{byzantine-mentorship}

These mechanisms inherit the BFT safety guarantees: with fewer than n/3
Byzantine committee members, the committed routing state is correct and
adapters derived from it are correctly generated. A residual risk
remains: a Byzantine coalition controlling fewer than n/3 seats could
generate subtly degrading adapters that pass fraud-proof verification
but introduce long-term drift. Detecting this requires longitudinal
performance monitoring across many checkpoints, a capability the
immutable checkpoint chain enables but whose analysis tooling is not yet
built. We flag this as an open problem rather than claiming it solved.

\subsection{5. Emergent Specialization}\label{emergent-specialization}

\textbf{{[}Hypothesis{]}} We hypothesize the Mentorship Protocol drives
emergent specialization through a reinforcing loop: nodes that perform
well on specific input classes receive more queries for those classes
(via routing), generate more training signal there, and accumulate more
specialized adapters. Over successive cycles this produces a network of
complementary experts rather than homogeneous generalists. This is the
organizational analogue of what Senge {[}3{]} calls team learning,
capabilities exceeding the sum of individual ones, with the routing
model as coordination mechanism ensuring specializations are
complementary rather than redundant. The paraconsistent aggregation of
Paper II preserves the information this needs: the Belnap state vector
tells the routing model where nodes agree (T), disagree (B), are
uncertain (N), or are confidently wrong (F). Whether emergent
specialization actually occurs, and whether it beats uniform model
replication, is an empirical question with methodology in Section 7.

\subsection{6. Relationship to the Gradient Papers
Series}\label{relationship-to-the-gradient-papers-series}

Paper I provides the infrastructure (GhostDAG, BFT checkpoints, the LVM
with AI precompiles, the adapter registry); the Mentorship Protocol
depends on it without modifying it. Paper II implements the protocol at
the protocol level: the Central Oracle becomes the BFT checkpoint,
mentor-mentee dynamics become the LoRA-adapter loop, and this paper
supplies the rationale for why a learned routing model (mentor) should
outperform a stateless aggregator (FedAvg). Paper VII (Mozi Cooperative)
extends the knowledge-sharing framework to economics: the principle that
value is created through knowledge transfer, not hoarding, is the
organizational foundation for cooperative ownership. Paper IX (Medusa
Paradigm) provides the biological inspiration: the jellyfish nerve net's
combination of local autonomy and global coordination through periodic
synchronization pulses is the biological analogue of the
checkpoint-based mentorship cycle, an inspirational analogy, not a
formal justification.

\subsection{7. Experimental Hypotheses}\label{experimental-hypotheses}

Earlier versions of this paper reported simulation figures (faster task
completion, error reduction, faster adaptation) that were placeholders
and did not represent measured outcomes. We keep them removed and state
three explicitly labelled hypotheses with designs.

\textbf{H1 --- Targeted mentorship outperforms flat aggregation.}
\emph{Claim:} a mentorship-driven swarm converges faster to a given
accuracy threshold than flat FedAvg on a heterogeneous task
distribution. \emph{Rationale:} targeted guidance reduces the search
space for improvement {[}3, 5{]}; a targeted adapter provides a gradient
direction pre-computed from the mentor's expertise, where FedAvg
requires the mentee to discover it. \emph{Methodology:} N = 50 non-IID
nodes on the testnet, comparing FedAvg baseline and the Mentorship
Protocol; measure rounds to a 90\% accuracy threshold, final accuracy,
per-domain variance, and total communication overhead, over 10 runs with
confidence intervals. \emph{Expected:} fewer rounds but higher per-round
overhead; whether total cost is lower is a prediction, not a result.

\textbf{H2 --- Double-loop learning improves routing over time.}
\emph{Claim:} a trainable routing model beats a fixed routing function
after sufficient training. \emph{Methodology:} compare static
(round-robin), fixed-weight (blue-score proportional), and learned
routing over \textasciitilde10,000 checkpoints on routing accuracy; if
learned routing does not beat fixed-weight routing significantly, the
double-loop mechanism adds complexity without benefit.

\textbf{H3 --- Emergent specialization produces complementary experts.}
\emph{Claim:} over successive cycles, mentorship produces greater
specialization diversity than a network without it. \emph{Methodology:}
measure specialization via the Herfindahl index of per-node accuracy
profiles, compare mentorship against flat aggregation over
\textasciitilde10,000 checkpoints, and test whether specialization is
complementary (different classes) or redundant (same classes);
complementary specialization is the predicted outcome, redundant would
indicate a failure of the routing model's diversity regularization.

\subsection{8. Conclusion}\label{conclusion}

The Mentorship Protocol provides the organizational-design rationale for
the Paraconsistent Consensus architecture of Paper II. Its central
contribution, the Central Oracle realized as a BFT finality checkpoint,
transforms stateless federated aggregation into a stateful, adaptive,
knowledge- accumulating system. The mappings from Senge, Nonaka and
Takeuchi, and Argyris and Schön provide the theoretical frame: systems
thinking explains why modeling interactions matters, the SECI spiral how
knowledge transforms between tacit and explicit, and double-loop
learning why the routing model must be trainable. We have been explicit
about the analogies' boundaries, strong for Combination and
Internalization, weak for Socialization, and about status: the mechanism
is now built (mentor loop, adapter lifecycle, on-chain registry), and
the organizational-outcome predictions are hypotheses, not results. The
fabricated metrics of earlier versions stay removed. What remains is the
empirical validation the infrastructure now makes possible.

\subsection{Acknowledgments}\label{acknowledgments}

Drafting and literature triage were assisted by AI systems; all
mechanical claims were verified by the authors against the referenced
source files and deployed contracts on chain 40204. This work received
no external funding.

\subsection{References}\label{references}

{[}1{]} Klosowski, L. (2023). Mentor/Mentee Relativity: Organizational
Learning in Mentorship-Driven Swarms. Unpublished internal note, Citrate
Inc. {[}2{]} McMahan, B., et al.~(2017). Communication-efficient
learning of deep networks from decentralized data. \emph{AISTATS},
1273--1282. {[}3{]} Senge, P. M. (1990). \emph{The Fifth Discipline: The
Art and Practice of the Learning Organization.} Doubleday. {[}4{]}
Nonaka, I., \& Takeuchi, H. (1995). \emph{The Knowledge-Creating
Company.} Oxford University Press. {[}5{]} Argyris, C., \& Schön, D. A.
(1978). \emph{Organizational Learning: A Theory of Action Perspective.}
Addison-Wesley. {[}6{]} da Silva, L. P. (2005). A formal model for the
Fifth Discipline. \emph{Journal of Artificial Societies and Social
Simulation (JASSS)}, 8(3), 6. \emph{(Verify locator before final
submission.)} {[}7{]} Klosowski, L., Mendenhall, L. (2026). Citrate:
Protocol Specification for an AI-Native BlockDAG Network. \emph{The
Gradient Papers No.~I} (this series). {[}8{]} Klosowski, L., Mendenhall,
L. (2026). Paraconsistent Consensus. \emph{The Gradient Papers No.~II}
(this series). {[}9{]} Hu, E. J., et al.~(2021). LoRA: Low-Rank
Adaptation of Large Language Models. \emph{ICLR 2022.} {[}10{]}
Blanchard, P., et al.~(2017). Machine learning with adversaries:
Byzantine-tolerant gradient descent. \emph{NeurIPS.} {[}11{]} Demers,
A., et al.~(1987). Epidemic algorithms for replicated database
maintenance. \emph{ACM PODC}, 1--12. {[}12{]} Shazeer, N., et
al.~(2017). Outrageously large neural networks: the sparsely-gated
mixture-of-experts layer. \emph{ICLR.} {[}13{]} Sompolinsky, Y., \&
Zohar, A. (2018). PHANTOM and GHOSTDAG. \emph{IACR ePrint.} {[}14{]}
Kaspa Network. (2025). KIP-14: Crescendo Hardfork. Activated May 5,
2025. {[}15{]} Belnap, N. D. (1977). A useful four-valued logic. In
Dunn, J. M., Epstein, G. (eds), \emph{Modern Uses of Multiple-Valued
Logic}, Reidel. {[}16{]} Biderman, D., et al.~(2024). LoRA learns less
and forgets less. \emph{Transactions on Machine Learning Research.}
{[}17{]} Ilharco, G., et al.~(2023). Editing models with task
arithmetic. \emph{ICLR.} {[}18{]} Bonabeau, E., Dorigo, M., Theraulaz,
G. (1999). \emph{Swarm Intelligence: From Natural to Artificial
Systems.} Oxford University Press. {[}19{]} Busoniu, L., Babuska, R., De
Schutter, B. (2008). A comprehensive survey of multiagent reinforcement
learning. \emph{IEEE Trans. Syst. Man Cybern.} 38(2), 156--172. {[}20{]}
Kirkpatrick, J., et al.~(2017). Overcoming catastrophic forgetting in
neural networks. \emph{PNAS}, 114(13), 3521--3526.

\begin{center}\rule{0.5\linewidth}{0.5pt}\end{center}

\emph{This paper is part of the Gradient Papers series, published by
Citrate Inc.} \emph{Correspondence: Larry@citrate.ai}
